\documentclass{beamer}

\input{settings.tex}


\title{Control Theory: recap}
\subtitle{Advanced Control}
\author{by Sergei Savin}
\centering
\date{\mydate}



\begin{document}
\maketitle

\begin{frame}{Content}
	\begin{itemize}
		\item Plants and Dynamical Systems
		\item Autonomous and Controlled Plants
		\item Controllers and State Feedback
		\item Closed-Loop Systems
		\item Motivation for Controller Design and Asymptotic Stability
		\item Linear Time-Invariant (LTI) Systems
		\item Stability of Linear Systems
		\item Stabilizing Control Design
		\item Gain Tuning Methods: Pole Placement and LQR
		\item Linearization of Nonlinear Systems
		\item Limitations of Linearization-Based Control
	\end{itemize}
\end{frame}


\begin{frame}{Plants and Dynamical Systems}
	\begin{flushleft}
		
		A \textbf{plant} is a physical system whose behaviour evolves over time.
		We describe it using a differential equation
		\[
		\dot{x}=f(x,u),
		\]
		where \(x\) is the \textbf{state} of the plant and \(u\) is its \textbf{input}.
		
		\bigskip
		
		{\mycolor{mydarkblue}
			
		\textbf{Example: Serial RC circuit}
		
		Let \(x\) be the voltage across the capacitor and \(u\) the input voltage.
		Then
		\[
		\dot{x}=-\frac{1}{RC}x+\frac{1}{RC}u.
		\]
	}
	\end{flushleft}
\end{frame}


\begin{frame}{Autonomous and Controlled Plants}
	\begin{flushleft}
		
		A plant described by
		\[
		\dot{x}=f(x)
		\]
		has no external control input. Such a system is called \textbf{autonomous}.
		
		\bigskip
		
		A plant described by \(\dot{x}=f(x,u)\) has an input \(u\) that can influence
		its behaviour. Behavious of an autonomous system on the other hand, is determined by its current state only.
		
		\bigskip
		
		{\mycolor{mydarkblue}
			
		\textbf{Examples:}
		
		\begin{itemize}
			\item  A DC motor is a plant with an input: we can apply a voltage to control it.
			\item A free pendulum is autonomous: its motion evolves according to its current state, without any control input.
		\end{itemize}
	}
		
	\end{flushleft}
\end{frame}


\begin{frame}{Controllers}
	\begin{flushleft}
		
		A \textbf{controller} is a rule that determines which input to apply to a plant.
		
		For a state-feedback controller, we write
		\[
		u=g(x).
		\]
		In other words, the controller observes the current state \(x\) and decides
		what control input \(u\) to apply.
		
		\bigskip
		
		{\mycolor{mydarkblue}
		\textbf{Example: Position control of a DC motor}
		
		A proportional controller
		\[
		u=-k_p(\varphi-\varphi^\star)
		\]
		applies a voltage proportional to the position error.
	}
	\end{flushleft}
\end{frame}


\begin{frame}{Closed-Loop Systems}
	\begin{flushleft}
		
		Consider a plant with an input:
		\[
		\dot{x}=f(x,u).
		\]
		If we connect it to a controller \(u=g(x)\), then
		\[
		\dot{x}=f\bigl(x,g(x)\bigr).
		\]
		
		The resulting system is called the \textbf{closed-loop system}. Since the
		input is now determined by the state, the closed-loop system is autonomous.
		
		\bigskip
		
		{\mycolor{mydarkblue}
		\textbf{Example}
		
		A water heater with built-in temperature regulation continuously adjusts its
		heating input based on the current temperature - autonomous closed-loop system.
	}
	
	\end{flushleft}
\end{frame}


\begin{frame}{Why Do We Design Controllers?}
	\begin{flushleft}
		
		One of the first goals of control design is to make the closed-loop system
		behave in a desirable way. In particular, we often want an equilibrium point \(x^\star\) to be
		\textbf{asymptotically stable}.
		
		\begin{block}{Simplified definition}
			If every solution starting sufficiently close to \(x^\star\) converges to
			\(x^\star\), then the equilibrium is asymptotically stable.
		\end{block}
		
		\bigskip
		
		{\mycolor{mydarkblue}
		\textbf{Example}
		
		An aircraft that returns to its desired orientation after a small perturbation
		is an example of stable closed-loop behaviour.
	}
		
	\end{flushleft}
\end{frame}


\begin{frame}{Linear Systems}
	\begin{flushleft}
		
		A particularly important class of plants is \textbf{linear time-invariant (LTI)}
		systems:
		\[
		\dot{x}=Ax+Bu.
		\]
		
		A common choice of controller is a \textbf{linear state-feedback law}
		\[
		u=-Kx.
		\]
		
		Substituting the controller into the plant gives the closed-loop system
		\[
		\dot{x}=(A-BK)x.
		\]
		
	\end{flushleft}
\end{frame}

\begin{frame}{Stability of Linear Systems}
	\begin{flushleft}
		
		Consider an autonomous LTI system
		\[
		\dot{x}=Ax.
		\]
		
		Its equilibrium \(x^\star=0\) is \textbf{asymptotically stable} if and only if
		all eigenvalues of \(A\) have strictly negative real parts:
		\[
		\Re\{\lambda_i(A)\}<0,
		\qquad \forall i.
		\]
		
		If the system is asymptotically stable, its solutions converge to zero
		\textbf{exponentially fast}.
		
	\end{flushleft}
\end{frame}

\begin{frame}{Stabilizing Control Design}
	\begin{flushleft}
		For the linear plant
		\[
		\dot{x}=Ax+Bu,
		\]
		and the controller \(u=-Kx\), the closed-loop system is
		\[
		\dot{x}=(A-BK)x.
		\]
		
		The stabilizing control problem is therefore:
			%
		\[
		\hspace{-0.5cm} \boxed{\text{Find }K\text{ such that all eigenvalues of }A-BK
			\text{ have negative real parts.}}
		\]
		
		We can solve this problem numerically using many algorithms, for example
		\textbf{pole placement} or \textbf{LQR}.
		
	\end{flushleft}
\end{frame}

\begin{frame}{Gain Tuning Methods}
	\begin{flushleft}
		
		Two common methods for finding a stabilizing feedback gain \(K\) are
		pole placement and LQR.
		
		\bigskip
		
		\textbf{Pole placement}:
		
		Finds \(K\) such that the closed-loop matrix \(A-BK\) has exactly the
		desired eigenvalues.
		
		\bigskip
		
		\textbf{Linear Quadratic Regulator (LQR)}:
		
		Finds \(K\) that minimizes the cost
		\[
		J=\int_0^\infty \left(x(t)^\top Qx(t)+u(t)^\top Ru(t)\right)dt.
		\]
		
	\end{flushleft}
\end{frame}

\begin{frame}{Linearization}
	\begin{flushleft}
		
		Most physical systems that we want to control are \textbf{nonlinear}.
		
		To apply linear control methods, we often approximate the nonlinear
		dynamics near an operating point by a linear system.
		
		Consider
		\[
		\dot{x}=f(x,u),
		\]
		with an equilibrium \((x^\star,u^\star)\).
		
		The linearized dynamics for small deviations are
		\[
		\dot{x}=Ax+Bu,
		\qquad
		A=\left.\frac{\partial f}{\partial x}\right|_{(x^\star,u^\star)},
		\quad
		B=\left.\frac{\partial f}{\partial u}\right|_{(x^\star,u^\star)}
		\]
		
	\end{flushleft}
\end{frame}

\begin{frame}{Limitations of Linearization-Based Control}
	\begin{flushleft}
		
		Consider the nonlinear plant
		\[
		\dot{x}=x+x^3+u.
		\]
		
		Linearizing around \(x^\star=0\) gives
		\[
		\dot{x}=x+u.
		\]
		
		Using pole placement, we can choose
		\[
		u=-2x,
		\]
		\hspace{-0.7cm} which gives the stable linear closed-loop system with eigenvalue \(-1\):
		\[
		\dot{x}=-x,
		\qquad \lambda=-1.
		\]
		
		But the closed-loop form for the original nonlinear system,
		\[
		\dot{x}=-x+x^3.
		\]
		If \(x(0)>1\), then \(\dot{x}>0\), and the solution moves away from zero.
		
	\end{flushleft}
\end{frame}






\begin{frame}{Literature}

\begin{itemize}
	
	
	\item MIT OpenCourseWare -- Feedback Control Systems: \bref{https://ocw.mit.edu/courses/16-30-feedback-control-systems-fall-2010/resources/mit16_30f10_lec05/}
	{Lecture 5: State-Space Systems}
	
	\item MIT OpenCourseWare -- Feedback Control Systems: \bref{https://ocw.mit.edu/courses/16-30-feedback-control-systems-fall-2010/resources/mit16_30f10_lec18/}
	{Lecture 18: Deterministic Linear Quadratic Regulator (LQR).}
	
	\item MIT OpenCourseWare -- Underactuated Robotics: \bref{https://ocw.mit.edu/courses/6-832-underactuated-robotics-spring-2009/72bc06c4dc73315bf49c28a81dc2b996_MIT6_832s09_read_ch03.pdf}
	{Chapter 3: Linear Quadratic Regulator, with discussion of linearized dynamics and LQR state feedback.}
	
\end{itemize}


\end{frame}



\myqrframe





\end{document}
